% !TeX program = lualatex
% =============================================================================
% TEMPLATE: IEEE-Style Two-Column CJK Technical Document (LuaLaTeX + babel)
%
% ENGINE: LuaLaTeX, not XeLaTeX.  babel's `onchar=ids fonts` — the mechanism
% that switches fonts per character between CJK and Latin — is documented as
% LuaTeX-only.  Under XeLaTeX the option is silently inert: the main language's
% font takes the whole document and no per-script switching happens.
% See references/troubleshooting.md, entry E1, for the XeLaTeX-only fallback.
%
% Copy everything above the TITLE & AUTHORS section marker into your document.
% The block between that marker and "END PREAMBLE" is the example title, author and date:
% write your own there, plus a second \hypersetup{pdftitle=..., pdfauthor=...}
% (later keys override the placeholders set further up).
% =============================================================================

\documentclass[10pt, a4paper, twocolumn]{article}
\usepackage[a4paper, top=2cm, bottom=2cm, left=1.5cm, right=1.5cm]{geometry}
\setlength{\columnsep}{16pt}   % article default 10pt is too tight at this width

\usepackage{fontspec}

% --- LANGUAGES ---------------------------------------------------------------
% English is the MAIN language, so \figurename/\tablename/\refname/\today stay
% English ("Figure", "Table", "References") as IEEE requires.  Traditional
% Chinese is a secondary locale reached automatically through onchar.
%
% Do NOT use the `japanese` locale for Chinese: babel has chinese-traditional
% (zh-Hant) and chinese-simplified (zh-Hans).  The japanese locale would set
% captions to 図/表/参考文献 and dates to Japanese format.
%
% For Simplified Chinese: swap chinese-traditional -> chinese-simplified.
% For Japanese:           swap chinese-traditional -> japanese.
% For Korean:             swap chinese-traditional -> korean  (also needs a KR
%                         font below; Korean line-breaks on spaces, not chars).
% intraspace widens the glue babel sets between Han characters.  The default
% stretch is too small once a narrow two-column page also carries unbreakable
% \texttt{} identifiers: on a 12-page zh-Hant paper, underfull \hbox warnings at
% badness 10000 fell from 26 to 5 with this one change.  Keep the base at 0 so
% normal lines are untouched; only the stretch and shrink move.
\usepackage[english, provide=*]{babel}
\babelprovide[import, onchar=ids fonts,
              intraspace={0 .22 .10}]{chinese-traditional}

% Ellipses.  onchar assigns a character to a locale by its script, and both
% ellipsis code points are script "Common", so they stay in the Latin font:
% Noto Serif has no U+22EF at all (the glyph is silently dropped) and sets
% U+2026 as low Latin dots.  Route both to the Chinese locale so 「⋯⋯」 and
% 「……」 come from the CJK font, centred, as Chinese typography expects.
% In English sentences write \ldots instead of typing the character.
\babelcharproperty{"2026}{locale}{chinese-traditional}
\babelcharproperty{"22EF}{locale}{chinese-traditional}

% --- FONT SELECTION WITH FALLBACK -------------------------------------------
% Font family names differ by distribution channel.  The same typeface ships as
% "Noto Serif CJK TC" (noto-cjk repo, Linux packages, Overleaf) and as
% "Noto Serif TC" (Google Fonts, which is what a Windows user usually installs).
% Hard-coding one name makes fontspec abort with a FATAL "font cannot be found".
% \PickFont walks a list and takes the first family actually installed.
% \PickFont* is the quiet form for optional fonts: when nothing on the list is
% installed it writes an Info line to the log instead of a Warning.
\ExplSyntaxOn
\NewDocumentCommand \PickFont { s m m }
  {
    \tl_clear_new:N #2
    \clist_map_inline:nn {#3}
      { \tl_if_empty:NT #2 { \IfFontExistsTF {##1} { \tl_set:Nn #2 {##1} } { } } }
    \tl_if_empty:NT #2
      {
        \IfBooleanTF {#1}
          { \msg_info:nnn    { template } { nofont } {#3} }
          { \msg_warning:nnn { template } { nofont } {#3} }
      }
  }
\msg_new:nnn { template } { nofont }
  { None~of~these~fonts~is~installed:~#1~--~run~scripts/check-fonts~. }
\ExplSyntaxOff

\PickFont \LatinSerif {Noto Serif, TeX Gyre Termes, Times New Roman, DejaVu Serif}
\PickFont \LatinSans  {Noto Sans, TeX Gyre Heros, Arial, DejaVu Sans}
\PickFont \LatinMono  {Noto Sans Mono, DejaVu Sans Mono, Consolas, Courier New}
\PickFont \CJKSerif   {Noto Serif CJK TC, Noto Serif TC, Source Han Serif TC,
                       PMingLiU, MingLiU, Microsoft JhengHei}
\PickFont \CJKSans    {Noto Sans CJK TC, Noto Sans TC, Source Han Sans TC,
                       Microsoft JhengHei, PMingLiU}
\PickFont \CJKMono    {Noto Sans Mono CJK TC, Noto Sans TC, Microsoft JhengHei,
                       MingLiU}

% THE WEIGHT-AXIS TRAP.  Google Fonts ships Noto TC as ONE variable file
% (NotoSerifTC-VF.ttf) whose first named instance is styled "ExtraLight,Regular"
% at weight 40.  Asking for the family by name therefore lands on ExtraLight:
% Chinese body text renders visibly thinner than the Latin beside it, with no
% warning at all.  Rendering the same text grey and measuring ink coverage gives
% roughly 1.0 : 1.3 : 1.8 for default : wght=400 : wght=700 (absolute numbers
% depend on point size and DPI; the ratio is the signal).  The cleanest proof:
% an explicit wght=200 render is BYTE-IDENTICAL to the unpinned one, and 200 is
% the ExtraLight end of the axis.
%
% Do not try to check this with pdffonts.  Subset names come from the base
% font's PostScript name, so a correctly pinned run still reports
% "NotoSerifTC-ExtraLight".  Measure the rendered ink instead.
%
% Pinning the axis fixes it, gives a REAL bold instead of a smeared fake one,
% and is silently ignored by static families (verified on PMingLiU and
% Microsoft JhengHei), so the same line is safe whichever font \PickFont found.
% CJK families have no italic and no small caps, so all four faces plus sc are
% mapped back onto the same file: FakeSlant supplies IEEE's \itshape headings,
% and SmallCapsFont stops \scshape (used on every \section) from emitting
% "Font shape .../m/sc undefined".  Note the remaining visual caveat: in a
% section title mixing scripts, the Latin half small-caps and the Chinese half
% cannot — for mostly-Chinese headings, swap \scshape for \bfseries below.
%
% All three families must be declared for BOTH scripts.  Declaring only rm —
% the common mistake — leaves CJK inside \texttt, Verbatim blocks and
% \sffamily TikZ node labels in a Latin font with no CJK glyphs; those
% characters are then dropped from the PDF with only a log warning.
\babelfont{rm}[Ligatures=TeX]{\LatinSerif}
\babelfont{sf}[Ligatures=TeX]{\LatinSans}
\babelfont{tt}{\LatinMono}

\newcommand{\CJKaxis}[1]{
  UprightFeatures    = {RawFeature={axis={wght=400}}},
  BoldFont           = {*},
  BoldFeatures       = {RawFeature={axis={wght=#1}}},
  ItalicFont         = {*},
  ItalicFeatures     = {RawFeature={axis={wght=400}}, FakeSlant=0.2},
  BoldItalicFont     = {*},
  BoldItalicFeatures = {RawFeature={axis={wght=#1}}, FakeSlant=0.2},
  SmallCapsFont      = {*},
  SmallCapsFeatures  = {RawFeature={axis={wght=400}}}
}

% --- GLYPH FALLBACK (auto-detected) -----------------------------------------
% Regional subsets have holes.  Google's "Noto Serif TC" has no Japanese
% shinjitai 図 (U+56F3) and skips many simplified forms; a glyph the font lacks
% is dropped from the PDF with one "Missing character" log line and a clean
% exit.  luaotfload can hand such glyphs to a second font.  The candidates
% below cover both platforms: full-repertoire Noto CJK on Linux/Overleaf,
% Microsoft JhengHei or Yu Gothic on Windows.  If none is installed, nothing is
% declared and the template behaves exactly as before.
% Verified through babel's onchar: the fallback fires inside rm, sf and tt.
% Limit: a fallback glyph inside bold text comes out at the second font's
% regular weight.
\PickFont* \CJKFallback {Noto Sans CJK TC, Source Han Sans TC, Microsoft JhengHei,
                         Noto Sans CJK JP, Yu Gothic}
\expandafter\IfBlankTF\expandafter{\CJKFallback}
  {\newcommand{\CJKfallbackkeys}{}}
  {\directlua{luaotfload.add_fallback("cjkfallback", {"\CJKFallback:mode=harf;"})}%
   \newcommand{\CJKfallbackkeys}{, RawFeature={fallback=cjkfallback}}}

\babelfont[chinese-traditional]{rm}[\CJKaxis{700}\CJKfallbackkeys]{\CJKSerif}
\babelfont[chinese-traditional]{sf}[\CJKaxis{700}\CJKfallbackkeys]{\CJKSans}
\babelfont[chinese-traditional]{tt}[\CJKaxis{700}\CJKfallbackkeys]{\CJKMono}

% --- CORE PACKAGES -----------------------------------------------------------
\usepackage{amsmath}
\usepackage{booktabs}
\usepackage{tabularx}
\usepackage{siunitx}      % S columns: decimal-aligned numeric benchmark data
\usepackage{xcolor}
\usepackage{graphicx}
\usepackage{placeins}     % \FloatBarrier
\usepackage{enumitem}
\usepackage{fancyhdr}
\usepackage{titlesec}
\usepackage{caption}
\usepackage{subcaption}
\usepackage{tikz}
\usetikzlibrary{positioning, arrows.meta, calc, fit, backgrounds}
\usepackage{fancyvrb}
\usepackage{balance}
\usepackage{microtype}    % protrusion; measurably fewer overfull lines in
                          % narrow justified columns
\usepackage{xurl}         % URL line-breaking in narrow columns
% hyperref goes near the END of the preamble, then bookmark.  Loading it before
% titlesec/placeins/fancyhdr (as the previous version did) is against hyperref's
% documented ordering and misplaces anchors on unnumbered sections.
\usepackage{hyperref}
\usepackage{bookmark}

% Ragged right in narrow tabularx cells; justified text in a 2 cm column
% produces huge inter-word gaps or overfull boxes.
\newcolumntype{Y}{>{\raggedright\arraybackslash}X}
\sisetup{detect-all, table-number-alignment = center}

\fvset{fontsize=\scriptsize}
\setlist[itemize]{label=--, leftmargin=*, nosep}
\setlist[enumerate]{leftmargin=*, nosep}

% --- IEEE NUMBERING ----------------------------------------------------------
% \thesection must be redefined, not just the titlesec label.  Otherwise
% \ref{sec:x} prints "2" while the heading shows "II.".
\renewcommand{\thesection}{\Roman{section}}
\renewcommand{\thesubsection}{\thesection-\Alph{subsection}}
\renewcommand{\thesubsubsection}{\thesubsection.\arabic{subsubsection}}
\renewcommand{\thetable}{\Roman{table}}
\setcounter{secnumdepth}{4}

% --- IEEE HEADING HIERARCHY --------------------------------------------------
% NOTE ON \scshape AND \itshape: CJK fonts have neither a small-caps nor an
% italic shape.  A Chinese section title under \scshape silently falls back to
% the upright shape, so a mixed title renders half small-capped and half plain.
% For a document whose headings are mostly Chinese, replace \scshape with
% \bfseries and \itshape with \bfseries below.
\titleformat{\section}[block]
  {\centering\normalsize\scshape}{\thesection.}{0.5em}{}
\titlespacing*{\section}{0pt}{1.2ex plus 0.5ex}{0.8ex plus 0.2ex}

\titleformat{\subsection}[block]
  {\normalsize\itshape}{\Alph{subsection}.}{0.5em}{}
\titlespacing*{\subsection}{0pt}{1ex plus 0.3ex}{0.5ex plus 0.1ex}

\titleformat{\subsubsection}[runin]
  {\normalsize\itshape}{\arabic{subsubsection})}{0.5em}{}[:]
\titlespacing*{\subsubsection}{\parindent}{0.8ex plus 0.2ex}{0.5em}

\titleformat{\paragraph}[runin]          % IEEE level 4: "a) label:"
  {\normalsize\itshape}{\alph{paragraph})}{0.5em}{}[:]
\titlespacing*{\paragraph}{2em}{0.6ex plus 0.2ex}{0.5em}

% --- IEEE CAPTIONS -----------------------------------------------------------
% Without this the article default renders "Figure 1: ..." with a colon at body
% size.  IEEE is "Fig. 1." below figures, and a centred "TABLE I" line with the
% title in small caps above tables.
\captionsetup[figure]{name={Fig.}, labelsep=period, font=footnotesize,
                      justification=raggedright, singlelinecheck=true}
\captionsetup[table]{labelsep=newline, font=footnotesize, textfont=sc,
                     justification=centering, name={TABLE}}
\captionsetup[subfigure]{labelformat=parens, labelsep=space, font=footnotesize}
\renewcommand{\thefigure}{\arabic{figure}}

% --- IEEE ABSTRACT & KEYWORDS ------------------------------------------------
% The environments emit the "Abstract—" and "Index Terms—" labels themselves.
% Do NOT type the label again in the body.
\renewenvironment{abstract}
  {\par\noindent\small\textbf{\textit{Abstract}---}\bfseries\ignorespaces}
  {\par\vspace{0.5em}}

\newenvironment{IEEEkeywords}
  {\par\noindent\small\textbf{\textit{Index Terms}---}\bfseries\ignorespaces}
  {\par\vspace{1em}\noindent\rule{\linewidth}{0.4pt}\par\vspace{0.5em}}

% Chinese abstract, for the 中英雙摘要 layout Taiwanese venues expect.
% Delete if the document is English-only.
\newenvironment{cnabstract}
  {\par\noindent\small\textbf{摘要---}\bfseries\ignorespaces}
  {\par\vspace{0.5em}}
\newenvironment{cnkeywords}
  {\par\noindent\small\textbf{關鍵詞---}\bfseries\ignorespaces}
  {\par\vspace{1em}}

% --- HEADER / FOOTER ---------------------------------------------------------
\pagestyle{fancy}
\fancyhf{}
\fancyfoot[C]{\small\thepage}
\renewcommand{\headrulewidth}{0pt}   % 0.4pt would draw a bare rule above an
                                     % empty header on every page
% To use a running header, uncomment both lines below AND set headrulewidth
% back to 0.4pt:
% \fancyhead[L]{\small Journal / Conference Name}
% \fancyhead[R]{\small Vol.~X, No.~Y, Month Year}

% --- FLOAT PLACEMENT TUNING (CRITICAL for two-column) ------------------------
\renewcommand{\topfraction}{0.9}
\renewcommand{\bottomfraction}{0.8}
\renewcommand{\textfraction}{0.1}
\renewcommand{\floatpagefraction}{0.8}
\setcounter{topnumber}{3}
\setcounter{bottomnumber}{2}
\setcounter{totalnumber}{5}

% --- HYPERLINKS & PDF METADATA -----------------------------------------------
% pdftitle must be plain text: it cannot be derived from \title here because
% \title contains \textbf and \\.
\hypersetup{
    colorlinks   = true,
    linkcolor    = blue!70!black,
    urlcolor     = blue!70!black,
    citecolor    = blue!70!black,
    pdftitle     = {Paper Title Here},
    pdfauthor    = {First A. Author; Second B. Author},
    pdfsubject   = {Technical report},
    pdfkeywords  = {keyword one, keyword two},
    bookmarksnumbered = true,
    bookmarksopen     = true,
}

% --- TITLE & AUTHORS ---------------------------------------------------------
% \thanks and footnotes are LOST inside the \twocolumn[\begin{@twocolumnfalse}]
% block below — the material is typeset in a box outside the main galley.  Put
% funding/corresponding-author notes in the affiliation lines, or emit them with
% \footnotetext after the block.
\title{\textbf{Paper Title Here: Subtitle if Needed\\
論文標題：如有副標題請加入}}
\author{
  First A. Author\\
  Department, University/Organization\\
  City, Country\\
  \texttt{first.author@example.com}
  \and
  Second B. Author\\
  Department, University/Organization\\
  City, Country\\
  \texttt{second.author@example.com}
}
\date{}   % IEEE papers carry no date under the title.  Comment out to show one.

% ============================== END PREAMBLE =================================
\begin{document}
\raggedbottom  % CRITICAL: without it LaTeX stretches inter-paragraph glue to
               % equalise column bottoms, gouging whitespace around floats.

\twocolumn[
  \begin{@twocolumnfalse}
    \maketitle
    \vspace{-1.5em}

    \begin{abstract}
    This paper presents\ldots{} (150--250 words). No citations, equations, or
    figure/table references belong in an abstract.
    \end{abstract}

    \begin{IEEEkeywords}
    keyword one, keyword two, keyword three
    \end{IEEEkeywords}

    % Delete the two blocks below for an English-only document.
    \begin{cnabstract}
    本文提出……（150--250 字）。摘要中不列引用、方程式或圖表參考。
    \end{cnabstract}

    \begin{cnkeywords}
    關鍵詞一、關鍵詞二、關鍵詞三
    \end{cnkeywords}
  \end{@twocolumnfalse}
]

% =============================================================================
% BODY
% =============================================================================

\section{Introduction}
\label{sec:intro}

這是一段中英混合文本範例。babel 依字元所屬文字系統自動切換中文與 English
字型，中英交界處請自行輸入一個半形空格。IEEE papers follow a standard
structure: Introduction, Related Work, Methodology, Results, Discussion,
Conclusion, References.

Cross-reference examples: as shown in Fig.~\ref{fig:architecture}, the
architecture\ldots{} The results in Table~\ref{tab:benchmark} show\ldots{}
The relationship is defined by \eqref{eq:example}.
Section~\ref{sec:method} now prints as ``Section III'' and matches its heading.

\section{Related Work}
\label{sec:related}

Prior studies~\cite{ref1} have explored\ldots{} Further work~\cite{ref2}
demonstrated improvements in\ldots{}

\subsection{Background}
This subsection uses the IEEE level-2 heading style.

\subsubsection{Specific Details}
A level-3 run-in heading; text continues after the colon.

\paragraph{Finer Points}
A level-4 run-in heading, IEEE's ``a)'' form.

\section{Methodology}
\label{sec:method}

The core relationship is
\begin{equation}
E = mc^2
\label{eq:example}
\end{equation}

Aligned equations:
\begin{align}
f(x) &= ax^2 + bx + c \label{eq:quad} \\
g(x) &= \frac{\mathrm{d}f}{\mathrm{d}x} = 2ax + b \label{eq:deriv}
\end{align}

CJK in math must be wrapped in \verb|\text{}|; babel's onchar switching acts on
text, not on math atoms: $P_{\text{尖峰}} = \max_t P(t)$.

\subsection{Implementation}

\begin{Verbatim}[fontsize=\scriptsize]
# docker-compose.yml — CJK 註解在此也能正常顯示
services:
  app:
    image: myapp:latest
    ports:
      - "8080:8080"
\end{Verbatim}

\begin{figure}[!t]
\centering
\resizebox{\columnwidth}{!}{%
\begin{tikzpicture}[
    >=Stealth, node distance=0.55cm,
    primary/.style={draw=blue!50!cyan!60, fill=blue!6, rounded corners=5pt,
        minimum width=4.2cm, minimum height=0.8cm,
        font=\sffamily\small, text=blue!60!black, line width=0.6pt, align=center},
    io/.style={draw=black!40, fill=white, rounded corners=6pt,
        minimum width=4.6cm, minimum height=0.72cm,
        font=\sffamily\small, text=black!80, line width=0.6pt},
    junction/.style={circle, draw=red!50, fill=red!5,
        inner sep=0pt, minimum size=0.5cm,
        font=\sffamily\bfseries\small, text=red!60!black},
    arr/.style={->, black!50, line width=0.6pt},
    skiparr/.style={-, black!30, line width=0.5pt},
]
\node[io] (input) {Input 輸入};
\node[primary, below=of input] (proc) {\textbf{Processing Block 處理區塊}};
\node[junction, below=of proc] (add) {\small +};
\node[io, below=of add] (output) {Output 輸出};

\begin{scope}[on background layer]
    \node[draw=black!30, dashed, rounded corners=8pt,
          inner xsep=0.4cm, inner ysep=0.35cm,
          fit=(proc)(add),
          label={[font=\sffamily\scriptsize\color{black!50},
                  anchor=north west]north west:{$\times$\,N layers}}
         ] (loopbox) {};
\end{scope}

\draw[arr] (input) -- (proc);
\draw[arr] (proc) -- (add);
\draw[arr] (add) -- (output);

\draw[skiparr]
    ($(input.south)!0.5!(proc.north)$)
    -| ([xshift=-0.25cm]loopbox.west |- add.west)
    -- (add.west);
\end{tikzpicture}%
}
\caption{Architecture diagram example. The processing block repeats $N$ times
with residual skip connections. CJK labels render because \texttt{sf} is
declared for the Chinese locale too.}
\label{fig:architecture}
\end{figure}

% --- SUBFIGURES: IEEE (a)/(b) panels sharing one "Fig. X." caption ---
\begin{figure}[!t]
\centering
\begin{subfigure}[t]{0.48\columnwidth}
  \centering
  \fbox{\rule{0pt}{1.6cm}\rule{0.8\linewidth}{0pt}}
  \caption{Before 處理前}
  \label{fig:panel-a}
\end{subfigure}\hfill
\begin{subfigure}[t]{0.48\columnwidth}
  \centering
  \fbox{\rule{0pt}{1.6cm}\rule{0.8\linewidth}{0pt}}
  \caption{After 處理後}
  \label{fig:panel-b}
\end{subfigure}
\caption{Paired panels. Reference a single panel as Fig.~\ref{fig:panel-a}.}
\label{fig:panels}
\end{figure}

\FloatBarrier
\section{Results}
\label{sec:results}

\begin{table}[!t]
\centering
\caption{Performance Benchmark Results}
\label{tab:benchmark}
\scriptsize
\setlength{\tabcolsep}{2.5pt}
\begin{tabularx}{\columnwidth}{@{}l S[table-format=3.1] S[table-format=1.2] Y@{}}
\toprule
\textbf{Method} & {\textbf{TG (t/s)}} & {\textbf{TTFT (s)}} & \textbf{Notes} \\
\midrule
\multicolumn{4}{@{}l}{\textit{A. Configuration Group}} \\[2pt]
Config A & 49.8 & 0.90 & Baseline 基準 \\
Config B\,\rlap{$^{\star}$}
         & 15.2 & 2.40 & \textbf{Recommended} 建議設定 \\
\bottomrule
\end{tabularx}

\vspace{6pt}
\raggedright
{\scriptsize
\textbf{TG} = token-generation throughput.\par\vspace{2pt}
$^{\star}$\,Default configuration for production use.
}
\end{table}

Table~\ref{tab:benchmark} shows Config~B gives the best throughput/latency
balance. Numeric columns are \texttt{siunitx} \texttt{S} columns, so the
decimal points align.

\section{Discussion}
\label{sec:discussion}

本節討論研究結果、局限性及其影響。

\section{Conclusion}
\label{sec:conclusion}

總結主要貢獻並建議未來研究方向。

\section*{Acknowledgment}
The authors thank\ldots{} 作者感謝……

% \appendix retitles sections as "Appendix A" only with the redefinition below;
% article's \appendix alone just switches \thesection to \Alph.
% \appendix
% \renewcommand{\thesection}{\Alph{section}}
% \titleformat{\section}[block]{\centering\normalsize\scshape}
%   {Appendix~\thesection.}{0.5em}{}
% \section{Supplementary Data}

% --- BIBLIOGRAPHY ------------------------------------------------------------
% thebibliography emits its own "References" heading — do not add \section*.
% \balance must appear in what will be the FIRST column of the LAST page; check
% the unbalanced output first and move it if the references start on the right.
\balance
\begingroup
\raggedright
\small
\begin{thebibliography}{99}
\bibitem{ref1}
A.~B. Author and C.~D. Author, ``Title of journal article,''
\textit{IEEE Trans. Example}, vol.~10, no.~3, pp.~100--110, Mar. 2024.

\bibitem{ref2}
E.~F. Author, ``Title of conference paper,'' in
\textit{Proc. IEEE Int. Conf. Example (ICE)}, City, Country, 2023, pp.~50--55.

\bibitem{ref3}
陳大文 and 林小明, ``中文期刊論文標題,'' \textit{中華民國資訊學會期刊},
vol.~5, no.~2, pp.~11--20, 2023 (in Chinese).

\bibitem{ref4}
I.~J. Author, ``Title of web resource,'' Website Name. [Online]. Available:
\url{https://example.com/resource}. [Accessed: Jan.~15, 2025].
\end{thebibliography}
\endgroup

\end{document}
